% Real slides: «Kryptologie» (Einführung), ported 1:1 from references/krypto_intro.odp.
% Slides that only reveal the next part of the same slide are one frame with overlays here.
% The original has no section slides; the two \section commands only feed the chrome.
% Pictures from the deck: theme/real/pictures/krypto/ (key, skytale). Polybios, Freimaurer and Caesar are drawn after the script; the punch card is typeset.
% Input by real-d.tex or real-b.tex, which define \RTitle, \RSection and \RTask; the drivers set the Kryptologie colour.
\slidesubject{Kryptologie}
\slidechapter{1}
\slidecoverimages{cover/kryptologie/binary-code.jpg, cover/kryptologie/enigma.jpg, cover/kryptologie/enigma-rotors.jpg,
  cover/kryptologie/padlock.jpg, cover/kryptologie/fingerprint.jpg, cover/kryptologie/vault-door.jpg}
\usetikzlibrary{shapes.symbols}
\title{Kryptologie}
\subtitle{Geheimschriften und Kryptosysteme}
\author{Florian Köhl}

% a box of the communication diagrams: \kbox{name}{position}{text}
\newcommand{\kbox}[3]{\node[fill=primaryDark, draw=borderColor, line width=0.6pt, rounded corners=\designInnerArc, text=white,
  minimum width=27mm, minimum height=8mm, font=\small] (#1) at (#2) {#3};}
% the communication diagram; #1: true for a cryptosystem (key, Verschlüsselung)
\newcommand{\kdiagram}[1]{%
  \begin{tikzpicture}[x=8.6mm, y=10mm]
    \node[font=\small, text=mutedFg] at (0,3.2) {Sender};
    \node[font=\small, text=mutedFg] at (9,3.2) {Empfänger};
    \kbox{k1}{0,2}{Klartext} \kbox{g1}{0,0}{Geheimtext}
    \kbox{k2}{9,2}{Klartext} \kbox{g2}{9,0}{Geheimtext}
    \draw[pointer, draw=borderColor] (k1) -- node[right, font=\footnotesize, text=textFg]
      {\ifx#1t\textbf{Ver}schlüsselung\else Chiffrierung\fi} (g1);
    \draw[pointer, draw=borderColor] (g2) -- node[left, font=\footnotesize, text=textFg]
      {\ifx#1t\textbf{Ent}schlüsselung\else Dechiffrierung\fi} (k2);
    \node[cloud, cloud puffs=11, draw=borderColor, fill=primaryDark, text=white, font=\small,
      minimum width=26mm, minimum height=11mm] (m) at (4.5,0) {Medium};
    \draw[edgeCol, line width=1pt, dashed, -{Stealth}] (g1) -- (m);
    \draw[edgeCol, line width=1pt, dashed, -{Stealth}] (m) -- (g2);
    \foreach \b in {k1,k2} \draw[borderColor, dashed, rounded corners=2mm]
      ([shift={(-4mm,-14mm)}]\b.south west |- g1.south) rectangle ([shift={(4mm,3mm)}]\b.north east);
    \ifx#1t
      \node[anchor=west, inner sep=0pt] at ([xshift=2mm]k1.east) {\includegraphics[height=9mm]{real/pictures/krypto/key.png}};
      \node[anchor=west, inner sep=0pt] at ([xshift=2mm]g2.east) {\includegraphics[height=9mm]{real/pictures/krypto/key.png}};
    \fi
  \end{tikzpicture}}

% Cipher figures after the script (scripts/Kryptologie: letter_rows.tex, caesar.sty): letters in typewriter
% bold on letter tiles; the theme's cell colours stand in for the script's tints, so they follow the palette.
\colorlet{kStrong}{nodeActive}   % coordinates, rims, dots: primaryDark in light, the term colour in dark palettes
\tikzset{k/tile/.style={cell, minimum size=9.6mm, inner sep=0pt, font=\ttfamily\bfseries\large, text=cellText},
  k/coord/.style={inner sep=0pt, font=\ttfamily\bfseries, text=kStrong},
  k/line/.style={draw=kStrong, line width=1.4pt, line cap=round},
  k/letter/.style={inner sep=0pt, font=\ttfamily\bfseries\large, text=textFg}}
% \PolybiosGrid: 5x5 square, row and column numbers outside (from 0, as in the script)
\newcommand{\PolybiosGrid}{%
  \begin{tikzpicture}[x=9.7mm, y=9.7mm]
    \foreach \row [count=\r from 0] in {{A,B,C,D,E},{F,G,H,IJ,K},{L,M,N,O,P},{Q,R,S,T,U},{V,W,X,Y,Z}}{%
      \node[k/coord] at (0,-\r) {\r};
      \foreach \L [count=\c] in \row \node[k/tile] at (\c,-\r) {\L};}
    \foreach \c in {0,...,4} \node[k/coord] at (\c+1,0.8) {\c};
  \end{tikzpicture}}
% \PigpenKey: the four figures of the Freimaurer alphabet in one row; a dot sits between letter and centre
\newcommand{\pigpendot}[1]{\fill[kStrong] (#1) circle (0.6mm);}
\newcommand{\PigpenKey}{%
  \begin{tikzpicture}[x=9.2mm, y=9.2mm]
    \foreach \x/\letters/\hasdot in {0/{A,B,C,D,E,F,G,H,I}/0, 3.5/{J,K,L,M,N,O,P,Q,R}/1}{%
      \draw[k/line] (\x+1,0) -- +(0,3) (\x+2,0) -- +(0,3) (\x,1) -- +(3,0) (\x,2) -- +(3,0);
      \foreach \L [count=\i from 0] in \letters {%
        \pgfmathsetmacro\cx{mod(\i,3)}\pgfmathsetmacro\cy{2-floor(\i/3)}
        \ifnum\hasdot=1
          \pgfmathsetmacro\dx{\cx==1 ? 0 : 1-\cx}\pgfmathsetmacro\dy{(\cx==1 && \cy==1) ? -1 : 1-\cy}
          \node[k/letter] at (\x+\cx+0.5-0.1*\dx,\cy+0.5-0.1*\dy) {\L};
          \pigpendot{\x+\cx+0.5+0.34*\dx,\cy+0.5+0.34*\dy}
        \else \node[k/letter] at (\x+\cx+0.5,\cy+0.5) {\L};\fi}}
    \foreach \x/\letters/\hasdot in {7.0/{S,T,U,V}/0, 10.5/{W,X,Y,Z}/1}{%
      \draw[k/line] (\x,0) -- +(3,3) (\x,3) -- +(3,-3);
      \foreach \L [count=\i] in \letters {%
        \ifcase\i\or\def\dx{0}\def\dy{1}\or\def\dx{-1}\def\dy{0}\or\def\dx{1}\def\dy{0}\or\def\dx{0}\def\dy{-1}\fi
        \node[k/letter] at (\x+1.5+0.9*\dx,1.5+0.9*\dy) {\L};
        \ifnum\hasdot=1 \pigpendot{\x+1.5+0.42*\dx,1.5+0.42*\dy}\fi}}
  \end{tikzpicture}}
% \CaesarWheel[k]: outer ring plaintext, inner ring ciphertext turned by k positions counter-clockwise;
% the pair at the top is marked
\NewDocumentCommand{\CaesarWheel}{O{3}}{%
  \begin{tikzpicture}[x=1.0cm, y=1.0cm]
    \def\R{2.12}\def\dR{0.58}\def\Ri{1.5}\def\dRi{0.58}
    \pgfmathsetmacro\step{360/26}
    \fill[cellFill] (0,0) circle (\R+\dR);
    \fill[cellActive] (0,0) circle (\Ri+\dRi);
    \fill[pageBg] (0,0) circle (\Ri);
    \foreach \a/\b in {\R/\dR, \Ri/\dRi}
      \fill[kStrong!50!pageBg] (90-\step/2:\a) arc[start angle=90-\step/2, end angle=90+\step/2, radius=\a]
        -- (90+\step/2:\a+\b) arc[start angle=90+\step/2, end angle=90-\step/2, radius=\a+\b] -- cycle;
    \foreach \L [count=\i from 0, evaluate=\i as \ang using 90-\i*\step] in {A,B,...,Z}{%
      \node[font=\ttfamily\bfseries, text=cellText, rotate=\ang-90] at (\ang:\R+\dR/2) {\L};
      \draw[cellRim, line width=0.5pt] (\ang+\step/2:\R) -- (\ang+\step/2:\R+\dR);}
    \foreach \L [count=\i from 0, evaluate=\i as \ang using 90-\i*\step+#1*\step] in {A,B,...,Z}{%
      \node[font=\ttfamily\bfseries, text=cellText, rotate=\ang-90] at (\ang:\Ri+\dRi/2) {\L};
      \draw[cellRim, line width=0.5pt] (\ang+\step/2:\Ri) -- (\ang+\step/2:\Ri+\dRi);}
    \draw[kStrong, line width=1pt] (0,0) circle (\R+\dR) (0,0) circle (\R) (0,0) circle (\Ri+\dRi) (0,0) circle (\Ri);
    \fill[kStrong] (0,0) circle (0.09);
  \end{tikzpicture}}

\begin{document}

\RTitle

\begin{frame}[text=large]{Ein bisschen Geschichte …}
  \begin{itemize}
    \only<1>{\item Wie alt ist die Kryptologie?}%
    \only<2>{\item Die \term{Kryptologie} ist (fast) so alt wie die Schrift.
    \begin{itemize}
      \item Erste Schriften ab ca. 2700 v.\,Chr.
      \item Lesen und Schreiben war lange ein Privileg.
      \item Mit zunehmender Bildung kam das Bedürfnis nach geheimer Kommunikation.
      \item Die Kryptologie war geboren!
    \end{itemize}}
  \end{itemize}
\end{frame}

\begin{frame}{Was ist Kryptologie?}
  \vfill
  \begin{center}
    \LARGE Kryptologie\\[2mm] $=$\\[2mm] \term{Kryptographie} $+$ \term{Krypt(o)analyse}
  \end{center}
  \vfill
\end{frame}

\begin{frame}[text=large]{Kryptographie}
  \begin{itemize}
    \item \term{«Verschleiern»} von Information
    \begin{itemize}
      \item Altgriechisch «geheim schreiben»
      \item Entwerfen von kryptographischen Systemen
      \item Umwandeln von Texten zu Geheimtexten
    \end{itemize}
  \end{itemize}
\end{frame}

\begin{frame}[text=large]{Krypt(o)analyse}
  \begin{itemize}
    \item \term{Information} aus «verschleierten» Texten \term{gewinnen}
    \begin{itemize}
      \item Analyse von kryptographischen Systemen
      \item Entschlüsseln von Geheimtexten
      \item Brechen (Knacken) von Kryptosystemen
    \end{itemize}
  \end{itemize}
\end{frame}

\begin{frame}[doit={zu zweit}, text=large]{Let's get cracking!}
  \begin{itemize}
    \item Was steht denn da?
    \begin{itemize}
      \item Versucht, die «geheimen» Botschaften zu knacken!
      \item Wie wurden diese erstellt?
    \end{itemize}
  \end{itemize}
\end{frame}

\section{Geheimschriften}

\begin{frame}[text=large]{Geheimschriften}
  \begin{itemize}
    \item Eine Geheimschrift ist ein \term{Paar von Algorithmen}.
    \begin{itemize}
      \item Algorithmus = «Verfahren»
    \end{itemize}
    \item \term{Chiffrierung}:
    \begin{itemize}
      \item Klartext → Geheimtext
    \end{itemize}
    \item \term{Dechiffrierung}:
    \begin{itemize}
      \item Geheimtext → Klartext
    \end{itemize}
  \end{itemize}
\end{frame}

\begin{frame}{Grundbegriffe}
  \vspace*{4mm}
  \centering\kdiagram{f}
\end{frame}

\begin{frame}[text=large]{Alphabete und Texte}
  \begin{itemize}
    \only<1>{\item Was ist eigentlich ein Text?}%
    \only<2>{%
    \item Ein \term{Alphabet} ist eine \textbf{Menge von Symbolen}.
    \begin{itemize}
      \item Lateinische Grossbuchstaben: $\{$A, B, C, …, X, Y, Z$\}$
      \item Das Computeralphabet: $\{0, 1\}$
    \end{itemize}
    \item Ein \term{Text} ist eine endliche \textbf{Folge von Symbolen} eines bestimmten Alphabets.
    \begin{itemize}
      \item \texttt{ICHBINEINTEXT}
      \item \texttt{01101001}
    \end{itemize}}
  \end{itemize}
\end{frame}

\begin{frame}{Geheimschriften}
  \begin{itemize}
    \item Eine Geheimschrift ist ein Paar von Verfahren / Algorithmen:
    \begin{itemize}
      \item \term{Chiffrierung}: Klartext → Geheimtext
      \item \term{Dechiffrierung}: Geheimtext → Klartext
    \end{itemize}
    \item Zusätzlich: Angabe des \term{Klartext}- und des \term{Geheimtextalphabets}
  \end{itemize}
\end{frame}

\begin{frame}[text=large]{Geheimschrift: Beispiel}
  \begin{itemize}
    \item Klartextalphabet: $\{$A, B, C, …, X, Y, Z$\}$
    \item Geheimtextalphabet:\\
      $\{$\foreach \c [count=\i] in {72,...,97} {\ding{\c}\ifnum\i<26 ,\,\allowbreak\fi}$\}$
    \item Chiffrierung / Dechiffrierung nach Position im Alphabet (A = \ding{72}, B = \ding{73}, …)
  \end{itemize}
\end{frame}

\begin{frame}[text=large]{Geheimschrift: Beispiel}
  \begin{itemize}
    \item \term{Atbash}
    \item Chiffrieren: A = Z, B = Y, C = X, …
    \item Dechiffrieren?
    \begin{itemize}
      \item \texttt{HKRVTOVRM HKRVTOVRM ZM WVI DZMW …}
    \end{itemize}
  \end{itemize}
\end{frame}

\begin{frame}[text=large]{Sicher?}
  \begin{itemize}
    \item Das Chiffrierungs- bzw. das Dechiffrierungs\term{verfahren} muss geheim gehalten werden.
    \item Wenn diese einem Angreifer bekannt sind, so ist die Geheimschrift nicht sicher!
  \end{itemize}
\end{frame}

\section{Kryptosysteme}

\begin{frame}[text=large]{Kerckhoffs' Prinzip}
  \begin{concept}[Kerckhoffs' Prinzip]
    Die Sicherheit eines Verschlüsselungsverfahrens darf \term{nicht} von der Geheimhaltung des Verfahrens
    selbst abhängen.
  \end{concept}
  \begin{itemize}
    \item Bei Geheimschriften ist dies nicht gegeben!
  \end{itemize}
\end{frame}

\begin{frame}[text=large]{Kryptosysteme}
  \begin{itemize}
    \item Bei einem \term{Kryptosystem} ist das Verschlüsselungsverfahren bekannt.
    \item Das Verfahren benötigt jedoch zusätzlich zum Klartext einen weiteren Input: \term{einen Schlüssel}.
    \item (Nur) der Schlüssel muss geheim gehalten werden!
  \end{itemize}
\end{frame}

\begin{frame}{Kryptosysteme}
  \vspace*{4mm}
  \centering\kdiagram{t}
\end{frame}

\begin{frame}{Beispiel: Lochkarten}
  \centering\vspace*{4mm}
  \begin{tikzpicture}[x=17mm, y=9mm]
    \foreach \row [count=\r] in {{A,T,X,O},{P,R,S,E},{I,G,V,M},{C,N,R,A},{E,T,L,O}}
      \foreach \l [count=\c] in \row
        \node[cell, minimum width=16mm, minimum height=8mm, font=\bfseries] (z\r\c) at (\c,-\r) {\l};
    % the punch card: everything except the holes T, O, P, S, E, C, R, E, T
    \begin{scope}[visible on=<2->]
      \foreach \r/\c in {1/1, 1/3, 2/2, 3/1, 3/2, 3/3, 3/4, 4/2, 4/4, 5/3, 5/4}
        \node[fill=mutedFg, draw=mutedFg, minimum width=16mm, minimum height=8mm, rounded corners=\designInnerArc] at (z\r\c) {};
    \end{scope}
  \end{tikzpicture}
\end{frame}

\begin{frame}[doit={Auftrag}, text=large]{Ihr seid dran!}
  \begin{itemize}
    \item Geheimschrift oder Kryptosystem?
    \item Gebt jeweils an:
    \begin{itemize}
      \item Ist es eine Geheimschrift oder ein Kryptosystem?
      \item Verfahren zur Verschlüsselung / Chiffrierung
      \item Bei KS: Was ist der Schlüssel?
    \end{itemize}
  \end{itemize}
\end{frame}

\begin{frame}{Geheimschrift der Freimaurer}
  \centering\vspace*{12mm}
  \PigpenKey
\end{frame}

\begin{frame}{Skytale}
  \centering\vspace*{2mm}
  \includegraphics[width=0.75\textwidth]{real/pictures/krypto/skytale.png}
\end{frame}

\begin{frame}{Polybios-Quadrat}
  \centering\vspace*{4mm}
  \PolybiosGrid
\end{frame}

\begin{frame}{Caesar-Verschlüsselung}
  \centering\vspace*{4mm}
  \CaesarWheel[3]
\end{frame}

\end{document}
